\chapter{1902: Emil Fischer}
\label{chap:chemistry1902}

The Nobel Prize in Chemistry for the year 1902 was awarded to Emil Fischer.

\textbf{Motivation:}
in recognition of the extraordinary services he has rendered by his work on sugar and purine syntheses

\section{Emil Fischer}

\subsection*{Early Life and Education}

Emil Fischer was born on 1852-10-09 in Euskirchen, Prussia (now Germany).
Hermann Emil Louis Fischer was a German chemist and 1902 recipient of the Nobel Prize in Chemistry. He co-developed the Fischer--Speier esterification with Arthur Speier in 1895. He also developed the Fischer projection, a symbolic way of drawing asymmetric carbon atoms, and proposed the lock-and-key model of enzyme action.

Fischer's early schooling was intended to prepare him for a business career, but his evident aptitude for science led him to study at the University of Bonn and then at the University of Strasbourg, where he worked under the organic chemist Adolf von Baeyer. It was in Strasbourg that Fischer received his training in the methods of organic synthesis and earned his doctorate in 1874. When Baeyer moved to Munich in 1875, Fischer followed him as an assistant, establishing himself as one of the most gifted younger chemists of the German school.

\subsection*{Scientific Context}

The last quarter of the nineteenth century saw organic chemistry transformed by two great achievements: the demonstration that carbon compounds could be synthesized from simple starting materials, and the establishment of the tetrahedral geometry of the carbon atom. August Kekulé's theory of the structure of benzene and van 't Hoff's stereochemical ideas had given chemists a new confidence that the architecture of even very complex natural substances could be worked out by systematic synthesis and degradation. At the same time, the sugar family and the related compounds known as purines were still largely uncharted, and the chemistry of proteins was almost entirely untouched. It was precisely these fields that Fischer made his own.

\subsection*{Career and Major Research}

Fischer held professorships at Erlangen, Würzburg, and, from 1892, Berlin. He determined the structures of the sugars, proving their stereochemistry and synthesizing glucose and other monosaccharides, and he worked out the chemistry of purines, synthesizing caffeine and related compounds. His investigations of amino acids and peptides, including the synthesis of peptides, made him a founder of carbohydrate and protein chemistry.

Fischer's work on sugars began in the 1880s and produced one of the great achievements of structural chemistry. He showed that the many sugars known at the time could be arranged into families related by the configuration of their asymmetric carbon atoms, and he determined the three-dimensional arrangement of those atoms by a combination of chemical degradation and synthesis. To represent these structures on paper, he introduced the convention now known as the Fischer projection, in which the carbon chain is drawn vertically and the position of each hydroxyl group indicates its spatial configuration. The stereochemistry of glucose and of the other common sugars was worked out in this way, and Fischer went on to synthesize glucose and several other sugars from simpler molecules.

In the chemistry of the purines, a group of nitrogen-containing ring compounds that includes uric acid and caffeine, Fischer carried out an equally systematic program. By building up the purine ring from simple starting materials and by relating the different members of the family to one another, he established the structure of uric acid, caffeine, theobromine, and related substances and synthesized many of them. This work created the chemistry of a whole class of natural products and eventually proved essential to the understanding of the nucleic acids, whose building blocks include purines.

\subsection*{Nobel Prize-winning Work}

The work for which Emil Fischer was awarded the Nobel Prize in Chemistry in 1902 was described by the Nobel Committee as: "in recognition of the extraordinary services he has rendered by his work on sugar and purine syntheses".

This research fundamentally advanced the field by elucidating the structures of sugars and purines through brilliant synthetic work, including the development of the Fischer projection for representing stereochemistry.

The synthesis of a sugar from smaller molecules was a feat of considerable difficulty because the chemist must not only join the atoms together but also control their spatial arrangement. Fischer's synthesis of glucose, announced in 1890, demonstrated that the methods of organic synthesis could reproduce the delicate stereochemistry of a natural product, and it crowned a decade of investigation in which he had first established the relationships among the sugars by careful analysis. The purine work was equally demanding: the synthesis of caffeine from simple starting materials required the construction of a fused ring system containing four nitrogen atoms, and it was carried out by a stepwise series of reactions that Fischer designed and executed with remarkable skill.

\subsection*{The Work in Detail}

Fischer's determination of sugar structures rested on two types of chemical evidence. The first was the study of the osazones, crystalline derivatives formed by sugars with phenylhydrazine, which Fischer and his collaborators used extensively in the 1880s. Because some pairs of sugars that differ only in the configuration of one carbon atom give the same osazone, these derivatives allowed Fischer to sort the sugars into related families and to prove which members of a family could be converted into one another. The second was the use of chemical degradation and synthesis to establish the configuration of the individual carbon atoms, culminating in the complete assignment of stereochemistry for the common sugars.

His work on peptides began at the turn of the century and was, if anything, even more ambitious. Fischer developed methods for joining amino acids together in chains by forming the amide linkage that joins them in proteins, and he succeeded in synthesizing peptides containing many amino acid residues. These compounds were the first synthetic substances bearing a real resemblance to proteins, and they allowed Fischer to study how the properties of such chains depend on their composition. His demonstration that a chain of more than a dozen amino acids could be built up step by step pointed the way toward the eventual synthesis of proteins themselves.

\subsection*{Reception}

Fischer's achievements were recognized throughout Europe with extraordinary speed. He was regarded by his contemporaries as the leading organic chemist of his generation, and his laboratory in Berlin attracted students from many countries. His syntheses of sugars and purines were admired not only as triumphs of skill but as models of how the chemistry of natural products ought to be pursued: by establishing structures through both synthesis and degradation, by preparing crystalline derivatives of known composition, and by relating every new substance to a general family. The Nobel Prize of 1902 confirmed a reputation that was already international.

\subsection*{Significance and Impact}

The impact of Emil Fischer's work extends far beyond the specific achievement recognized by the Nobel Prize.
Emil Fischer's elucidation of sugar structures and development of the Fischer projection became essential tools in organic chemistry and biochemistry. Emil Fischer's work on purines established the chemical basis for understanding nucleic acids.

Fischer's later research opened up entirely new territory. His studies of enzyme specificity in reactions involving sugars led him to formulate the lock-and-key model of enzyme action, in which the enzyme and its substrate fit together as a key fits a lock, a picture that remains a valuable first description of specificity in biological catalysis. His work on the chemistry of proteins, begun with the synthesis of peptides, laid the foundation of protein chemistry, and his demonstration that proteins are built from amino acids joined in amide linkages, joined with the later work of others, pointed the way to the modern understanding of these molecules.

\subsection*{Later Career and Honors}

Emil Fischer continued his scientific work until 1919-07-15, passing away in Berlin, Germany.
He received numerous honors including membership in major scientific academies and societies.
Emil Fischer served as president of the German Chemical Society and was instrumental in establishing the Kaiser Wilhelm Society.

From his chair in Berlin, which he occupied from 1892 until his death, Fischer directed one of the largest and most productive chemical institutes in the world. He was a central figure in the organization of German science, and he played a leading role in the founding of the Kaiser Wilhelm Society for the Advancement of Science, whose institutes carried on research in fields too large for the universities. The war years were difficult for him, as they were for German science generally, and he withdrew increasingly from public life before his death in 1919.

\subsection*{Legacy}

The legacy of Emil Fischer endures through the lasting impact of his scientific contributions.
The Fischer projection remains the standard method for depicting carbohydrate stereochemistry. Emil Fischer's work on enzyme specificity and the lock-and-key model laid the groundwork for modern biochemistry.

Fischer's influence can be traced through nearly every branch of organic and biological chemistry. His methods for determining the structure of natural products became the model for generations of chemists, and the fields he founded, carbohydrate chemistry, purine chemistry, and peptide chemistry, remain vigorous today. The nucleic acid chemistry that blossomed in the twentieth century rests directly on his elucidation of the purines, and the understanding of enzymes and proteins that is central to modern biology owes much to his early and prescient investigations.

\subsection*{Collaborators and Students}

Fischer trained a remarkable number of students who went on to distinguished careers, among them Hermann Leuchs, Otto Diels, and others who carried his methods into new areas. His insistence on rigorous purification, on the use of crystalline derivatives as proof of structure, and on the value of total synthesis as the ultimate test of an assignment, shaped the style of organic chemistry for half a century. Through these students, and through his many writings, Fischer's approach to the chemistry of natural products was transmitted to laboratories throughout the world.

\subsection*{Modern Developments}

The carbohydrate and purine chemistry advanced by Emil Fischer remains central to biochemistry and the pharmaceutical industry. Sugar chemistry underlies glycobiology, vaccine development, and the synthesis of antiviral drugs, while the Fischer--Speier esterification and Fischer's lock-and-key model of enzyme action are still taught as foundational concepts. Modern glycomics and synthetic carbohydrate vaccines build on the analytical framework Fischer established more than a century ago.

The synthesis of peptides that Fischer pioneered grew into the modern discipline of peptide chemistry, in which automated synthesizers build chains of amino acids for use as pharmaceuticals, research reagents, and materials. The Fischer projections used in every biochemistry textbook, the esterification reaction named after him, and the concepts of stereoisomerism that he worked out for the sugars are all in daily use. The lock-and-key image remains the first explanation offered to students of how enzymes achieve their remarkable selectivity, and the structural approach Fischer brought to biology remains one of the foundations of chemical biology.

\subsection*{Selected Publications}

"Ueber die Constitution des Caffeïns, Xanthins, Hypoxanthins und verwandter Basen" (1881)
"Synthesen in der Zuckergruppe" (1890)
"Ueber die Configuration des Traubenzuckers und seiner Isomeren" (1891)
"Untersuchungen über Aminosäuren, Polypeptide und Proteine" (1899--1906)
